% Source: physical PDF pages 81--85; printed pages 58--62.
% Problems begin after Appendix B on physical p. 81; References end on physical p. 85.
% Contents: Problems 1--9; References 1--26, including 11a and 13a (28 displayed entries).
% Figures/tables/footnotes: none. Source uncertainties: none.

\chapterbackmatter{Problems}

\begin{enumerate}
  \item Derive the Bianchi identity
  \[
    D_\mu F_{a\nu\lambda}+D_\nu F_{a\lambda\mu}
    +D_\lambda F_{a\mu\nu}=0.
  \]

  \item Suppose we use generalized Coulomb gauge in a non-Abelian gauge theory, taking the gauge-fixing function as $f_a=\boldsymbol\nabla\cdot\mathbf A_a$. Derive the ghost Lagrangian. What is the ghost propagator? (Take $B[f]=\exp(-\ii\int \dd^4x\,f_a f_a/2\xi)$.)

  \item Suppose that in electrodynamics we use a gauge-fixing function $f=\partial_\mu A^\mu+cA_\mu A^\mu$, with $c$ an arbitrary constant. Derive the ghost Lagrangian.\\
  (Take $B[f]=\exp(-\ii\int \dd^4x\,f_a f_a/2\xi)$.) What is the ghost propagator?

  \item Show that there is no simple Lie algebra with just four generators.

  \item Show that if a field $\psi_\ell(x)$ belonging to a representation of a gauge group with generator matrices $t_a$ varies along a path $x^\mu=x^\mu(\tau)$ according to the differential equation
  \[
    \frac{\dd\psi(\tau)}{\dd\tau}
    =\ii t_a\psi(\tau)A_{a\mu}(x(\tau))\frac{\dd x^\mu(\tau)}{\dd\tau},
  \]
  (with both $\psi_\ell(x)$ and $A_{a\mu}(x)$ classical c-number fields) then the change in $\psi$ around any small closed path $\gamma$ around a point $X^\mu$ is proportional to
  \[
    t_a\psi(X)F_{a\mu\nu}(X)
    \oint_c x^\mu(\tau)\frac{\dd x^\nu(\tau)}{\dd\tau}\,\dd\tau.
  \]
  Find the constant of proportionality. Use this result to show that if $F_{a\mu\nu}(x)$ vanishes everywhere then it is possible by a gauge transformation to make $A_{a\mu}(x)$ vanish throughout at least a finite region. (Hint: Follow the analogous argument in electrodynamics, or in general relativity, as in Ref.~\hyperref[ch15-ref-6]{[6]}.)

  \item Applying path-integral methods to a general non-Abelian gauge theory, calculate the propagator of the gauge field $A_{a\mu}(x)$ if we choose a gauge-fixing functional $f_a=n_\mu A_a^\mu$, with $n_\mu$ arbitrary constants. (Take $B[f]=\exp(-\ii\int \dd^4x\,f_a f_a/2\xi)$.) What is the ghost Lagrangian? What is the ghost propagator? What is the ghost interaction vertex?

  \item Suppose we adopt BRST invariance instead of gauge invariance as a fundamental physical principle. Derive the most general Lagrangian density constructed from sums of products of fields and field derivatives of dimensionality $(\text{mass})^d$ with $\dd\leq4$, constructed from $A_{a\mu}$, $\omega_a$, $\omega_a^*$, $h_a$, and/or their derivatives, that is invariant (up to total derivatives) under Lorentz transformations, ghost number phase transformations ($\omega_a\longrightarrow e^{\ii\theta}\omega_a$, $\omega_a^*\longrightarrow e^{-\ii\theta}\omega_a^*$), global gauge transformations ($\epsilon_a$ constant), and BRST invariance.

  \item Show that the antibracket satisfies the symmetry condition \eqref{eq:15.9.15} and the Jacobi identity \eqref{eq:15.9.21}.

  \item Show that if a functional $O$ satisfies the condition $(O,S)=\ii\Delta S$ and the action $S$ satisfies the quantum master equation then the quantum average $\langle O\rangle$ is independent of the gauge-fixing functional $\Psi$.
\end{enumerate}

\chapterbackmatter{References}

\begin{enumerate}
  \item \phantomsection\label{ch15-ref-1}
  C.~N.~Yang and R.~L.~Mills, \textit{Phys. Rev.} \textbf{96}, 191 (1954). O.~Klein had come close to a formulation of the $SU(2)$ Yang--Mills theory in a talk at a 1938 conference at Warsaw, reported in \textit{New Theories in Physics} (International Institute of Intellectual Cooperation, Paris, 1939). For a critical discussion of Klein's theory, see D.~J.~Gross, `Oscar Klein and Gauge Theory,' talk at the 1994 Oscar Klein symposium at Stockholm, Princeton preprint PUPT-1508.

  \item \phantomsection\label{ch15-ref-2}
  R.~Utiyama, \textit{Phys. Rev.} \textbf{101}, 1597 (1956).

  \item \phantomsection\label{ch15-ref-3}
  R.~P.~Feynman, \textit{Acta Phys. Polonica} \textbf{24}, 697 (1963).

  \item \phantomsection\label{ch15-ref-4}
  L.~D.~Faddeev and V.~N.~Popov, \textit{Phys. Lett.} \textbf{25B}, 29 (1967).

  \item \phantomsection\label{ch15-ref-5}
  B.~S.~De Witt, \textit{Phys. Rev.} \textbf{162}, 1195, 1239 (1967).

  \item \phantomsection\label{ch15-ref-6}
  The proof is easily adapted from the standard proof of the corresponding result in general relativity that, for the existence of a coordinate transformation to a flat metric in a finite simply connected region, it is necessary and sufficient that the Riemann--Christoffel curvature tensor should vanish. See, for example, S.~Weinberg, \textit{Gravitation and Cosmology} (Wiley, New York, 1972): Sections 6.3 and 6.4.

  \item \phantomsection\label{ch15-ref-7}
  The proof that \textbf{a} (with $g_{ab}=\delta_{ab}$) implies \textbf{b} and that \textbf{b} implies \textbf{c} was given by M.~Gell-Mann and S.~L.~Glashow, \textit{Ann. Phys. (N.Y.)} \textbf{15}, 437 (1961).

  \item \phantomsection\label{ch15-ref-8}
  See, for example, S.~Weinberg, \textit{ibid.}, Section 7.6.

  \item \phantomsection\label{ch15-ref-9}
  V.~Gribov, \textit{Nucl. Phys.} B\textbf{139}, 1 (1978). Also see R.~Jackiw, I.~Muzinich, and C.~Rebbi, \textit{Phys. Rev.} D\textbf{17}, 1576 (1978); R.~Jackiw, in \textit{New Frontiers in High Energy Physics}, eds. B.~Kursunoglu, A.~Perlmutter, and L.~Scott (Plenum, New York, 1978); N.~Christ and T.~D.~Lee, \textit{Phys. Rev.} D\textbf{22}, 939 (1980); R.~Jackiw, in \textit{Current Algebra and Anomalies} (World Scientific, Singapore, 1985): footnote 50.

  \item \phantomsection\label{ch15-ref-10}
  C.~Becchi, A.~Rouet, and R.~Stora, \textit{Comm. Math. Phys.} \textbf{42}, 127 (1975); in \textit{Renormalization Theory}, eds. G.~Velo and A.~S.~Wightman (Reidel, Dordrecht, 1976); \textit{Ann. Phys.} \textbf{98}, 287 (1976).

  \item \phantomsection\label{ch15-ref-11}
  I.~V.~Tyutin, Lebedev Institute preprint N39 (1975).

  \item[11a.] \phantomsection\label{ch15-ref-11a}
  N.~Nakanishi, \textit{Prog. Theor. Phys.} \textbf{35}, 1111 (1966); B.~Lautrup, \textit{Mat. Fys. Medd. Kon. Dan. Vid.-Sel. Medd.} \textbf{35}, 29 (1967).

  \item \phantomsection\label{ch15-ref-12}
  The original source of this argument is unknown to me. I learned it from J.~Polchinski.

  \item \phantomsection\label{ch15-ref-13}
  S.~N.~Gupta, \textit{Proc. Phys. Soc.} \textbf{63}, 681 (1950); \textbf{64}, 850 (1951); K.~Bleuler, \textit{Helv. Phys. Acta} \textbf{3}, 567 (1950); K.~Bleuler and W.~Heitler, \textit{Progr. Theor. Phys.} \textbf{5}, 600 (1950).

  \item[13a.] \phantomsection\label{ch15-ref-13a}
  G.~Curci and R.~Ferrari, \textit{Nuovo Cimento} \textbf{35}, 273 (1976); T.~Kugo and I.~Ojima, \textit{Prog. Theor. Phys.} \textbf{60}, 1869 (1978); \textit{Prog. Theor. Phys. Suppl.} \textbf{66}, 1 (1979). Also see Ref.~\hyperref[ch15-ref-11a]{[11a]} for the case of electrodynamics.

  \item \phantomsection\label{ch15-ref-14}
  J.~Thierry-Mieg, \textit{J. Math. Phys.} \textbf{21}, 2834 (1980); R.~Stora, in \textit{Progress in Gauge Theories: Proceedings of a Symposium at Carg\`ese, 1983}, eds. G.~'t Hooft, A.~Jaffe, H.~Lehmann, P.~K.~Mitter, I.~M.~Singer, and R.~Stora (Plenum, New York, 1984): p.~543; L.~Bonora and P.~Cotta-Ramusino, \textit{Commun. Math. Phys.} \textbf{87}, 589 (1983); J.~Ma\~nes, R.~Stora, and B.~Zumino, \textit{Commun. Math. Phys.} \textbf{102}, 157 (1985); A.~S.~Schwarz, \textit{Commun. Math. Phys.} \textbf{155}, 249 (1993); P.~M.~Lavrov, P.~Yu. Moshin and A.~A.~Reshetnyak, Tomsk preprint hep-th/9507104; P.~M.~Lavrov, Tomsk preprint hep-th/9507105.

  \item \phantomsection\label{ch15-ref-15}
  W.~Siegel, \textit{Phys. Lett.} \textbf{93B}, 170 (1980); T.~Kimura, \textit{Prog. Theor. Phys.} \textbf{64}, 357 (1980); \textbf{65}, 338 (1981).

  \item \phantomsection\label{ch15-ref-16}
  G.~Barnich, F.~Brandt, and M.~Henneaux, \textit{Commun. Math. Phys.} \textbf{174}, 57 (1995): Section 14.

  \item \phantomsection\label{ch15-ref-17}
  G.~Curci and R.~Ferrari, \textit{Nuovo Cimento} \textbf{32A}, 151 (1976); I.~Ojima, \textit{Prog. Theor. Phys. Supp.} \textbf{64}, 625 (1980); L.~Baulieu and J.~Thierry-Mieg, \textit{Nucl. Phys.} B\textbf{197}, 477 (1982).

  \item \phantomsection\label{ch15-ref-18}
  The anti-BRST transformation was extended to general local symmetries by L.~Alvarez-Gaum\'e and L.~Baulieu, \textit{Nucl. Phys.} B\textbf{212}, 255 (1983). It is induced by the operator:
  \begin{align*}
    \bar s={}&\omega^{*A}\delta_A\phi^r\frac{\delta_L}{\delta\phi^r}
    -\frac{1}{2}\omega^{*B}\omega^{*C}f^A{}_{BC}\frac{\delta_L}{\delta\omega^{*A}}
    -\omega^{*B}h^Cf^A{}_{BC}\frac{\delta}{\delta h^A}\\
    &+\left[-f^A{}_{BC}\omega^{*B}\omega^C+h^A\right]\frac{\delta_L}{\delta\omega^A}.
  \end{align*}
  It is straightforward to show that this transformation is nilpotent:
  \[
    \bar s^2=0.
  \]
  Also, the BRST and anti-BRST transformations anticommute. Furthermore, there is a cohomology theorem (unpublished) for the anti-BRST transformations, like that for the BRST transformations: the most general functional $I$ of $\phi^r$, $\omega^A$, $\omega^{*A}$, and $h^A$ that satisfies the anti-BRST invariance condition $\bar sI=0$ and has ghost number zero is of the form
  \[
    I[\phi,\omega,\omega^*,h]
    =I_0[\phi]+\bar s\bar\Psi[\phi,\omega,\omega^*,h].
  \]
  (Just use the Hodge operator $\bar t\equiv\omega^A\delta/\delta h^A$ in place of $t$.) Anti-BRST symmetry can be an aid in enumerating possible terms in Lagrangians and Greens functions, but I do not know of any case where it is indispensable.

  \item \phantomsection\label{ch15-ref-19}
  I.~A.~Batalin and G.~A.~Vilkovisky, \textit{Phys. Lett.} B\textbf{102}, 27 (1981); \textit{Nucl. Phys.} B\textbf{234}, 106 (1984); \textit{J. Math. Phys.} \textbf{26}, 172 (1985). Also see B.~L.~Voronov and I.~V.~Tyutin, \textit{Theor. Math. Phys.} \textbf{50}, 218, 628 (1982). For a lucid review, see J.~Gomis, J.~Par\'is and S.~Samuel, \textit{Phys. Rep.} \textbf{259}, 1 (1995).

  \item \phantomsection\label{ch15-ref-20}
  E.~S.~Fradkin and G.~A.~Vilkovisky, \textit{Phys. Lett.} B\textbf{55}, 224 (1975); CERN report TH2332 (1977); I.~A.~Batalin and G.~A.~Vilkovisky, \textit{Phys. Lett.} B\textbf{69}, 309 (1977); I.~S.~Fradkin and T.~E.~Fradkina, \textit{Phys. Lett.} B\textbf{72}, 334 (1977). Also see M.~Henneaux and C.~Teitelboim, \textit{Quantization of Gauge Systems} (Princeton University Press, Princeton, 1992).

  \item \phantomsection\label{ch15-ref-21}
  I.~A.~Batalin and I.~V.~Tyutin, \textit{Phys. Lett.} B\textbf{356}, 373 (1995).

  \item \phantomsection\label{ch15-ref-22}
  J.~Zinn-Justin, in \textit{Trends in Elementary Particle Theory---International Summer Institute on Theoretical Physics in Bonn 1974} (Springer-Verlag, Berlin, 1975): p.~2.

  \item \phantomsection\label{ch15-ref-23}
  D.~Z.~Freedman, P.~van Nieuwenhuizen, and S.~Ferrara, \textit{Phys. Rev.}, D\textbf{13}, 3214 (1976); R.~E.~Kallosh, \textit{Nucl. Phys.} B\textbf{141}, 141 (1978); B.~de Wit and J.~W.~van Holten, \textit{Phys. Lett.}, B\textbf{79}, 389 (1979); P.~van Nieuwenhuizen \textit{Phys. Rep.}, \textbf{68}, 189 (1981).

  \item \phantomsection\label{ch15-ref-24}
  M.~Henneaux and C.~Teitelboim, Ref.~\hyperref[ch15-ref-20]{[20]}: Section 18.1.4.

  \item \phantomsection\label{ch15-ref-25}
  E.~P.~Wigner, \textit{Group Theory} (Academic Press, New York, 1959): pp.~76--8.

  \item \phantomsection\label{ch15-ref-26}
  The original reference is E.~Cartan, \textit{Sur la Structure des Groupes de Transformations Finis et Continue}, (Paris, 1894; 2nd edn 1933). For a textbook treatment, see, for example, R.~Gilmore, \textit{Lie Groups, Lie Algebras, and Some of Their Applications} (Wiley, New York, 1974): Chapter 9.
\end{enumerate}
